
\chapter{Байєсівська фільтрація, згладжування і факторні графи}

\sourcebox{Джерельна лінія: огляд неперервно-часового оцінювання стану, Labbe Kalman book, Probabilistic Robotics-подібна нотація, факторні графи в сучасній робототехніці. Розділ є зведеним практичним містком, не повним перекладом одного джерела.}

\section{Фільтрація}

Фільтрація оцінює поточний стан за минулими і поточними вимірюваннями:
\[
p(x_k\mid z_{1:k},u_{1:k}).
\]
Байєсівський фільтр має два кроки.

\subsection{Прогноз}
\[
p(x_k\mid z_{1:k-1},u_{1:k})
=
\int p(x_k\mid x_{k-1},u_k)
p(x_{k-1}\mid z_{1:k-1},u_{1:k-1})\,\dd x_{k-1}.
\]

\subsection{Оновлення}
\[
p(x_k\mid z_{1:k},u_{1:k})
=
\eta\, p(z_k\mid x_k)
p(x_k\mid z_{1:k-1},u_{1:k}),
\]
де $\eta$ --- нормувальна константа.

Фільтр Калмана є спеціальним випадком, коли моделі лінійні, а розподіли гаусівські:
\begin{align}
x_k &= F_kx_{k-1}+B_ku_k+w_k,\qquad w_k\sim\Normal(0,Q_k),\\
z_k &= H_kx_k+r_k,\qquad r_k\sim\Normal(0,R_k).
\end{align}

\section{EKF як локально лінійний Байєсівський фільтр}

Для нелінійних моделей:
\[
x_k=f(x_{k-1},u_k,w_k),\qquad z_k=h(x_k,r_k).
\]
EKF замінює їх локальними лінеаризаціями:
\[
F_k = \left.\frac{\partial f}{\partial x}\right|_{\hat x_{k-1},u_k},\qquad
H_k = \left.\frac{\partial h}{\partial x}\right|_{\hat x_k^-}.
\]
У ESKF похідні беруть не за повним станом, а за похибкою:
\[
H_k = \left.\frac{\partial h(x\oplus\delta x)}{\partial\delta x}\right|_{\delta x=0}.
\]

\section{Згладжування}

Згладжування оцінює не лише поточний стан, а траєкторію або вікно траєкторії:
\[
p(x_{0:N}\mid z_{1:N},u_{1:N}).
\]
Це часто перетворюється на задачу оптимізації:
\[
x_{0:N}^{\star}
=
\arg\min_{x_{0:N}}
\left(
\sum_k \|r_k^{\mathrm{motion}}(x_{k-1},x_k,u_k)\|_{Q_k^{-1}}^2
+
\sum_k \|r_k^{\mathrm{meas}}(x_k,z_k)\|_{R_k^{-1}}^2
\right).
\]
Тут
\[
\|r\|_{\Sigma^{-1}}^2 = r\trans\Sigma^{-1}r.
\]
Фільтр придатний для низької затримки. Згладжування часто точніше і стабільніше, бо переоцінює попередні стани після приходу нових вимірювань.

\section{Факторний граф}

Факторний граф розбиває цільову функцію на локальні фактори:
\[
p(x\mid z)\propto \prod_i \phi_i(x_i).
\]
У логарифмічній формі:
\[
-\log p(x\mid z)=\sum_i \|r_i(x_i)\|_{\Sigma_i^{-1}}^2+\mathrm{const}.
\]
Після лінеаризації:
\[
r_i(x\oplus\delta x)\simeq r_i(x)+J_i\delta x.
\]
Нормальні рівняння:
\[
\left(\sum_i J_i\trans\Sigma_i^{-1}J_i\right)\delta x
=
-\sum_i J_i\trans\Sigma_i^{-1}r_i.
\]
На практиці цю систему розв'язують не явним оберненням, а факторизаціями QR/Cholesky або ітераційними методами.

\section{Фільтр проти згладжувача}

\begin{longtable}{L{0.24\linewidth}L{0.33\linewidth}L{0.33\linewidth}}
\toprule
Критерій & Фільтрація & Згладжування\\
\midrule
Стан & Поточний $x_k$ & Траєкторія або вікно $x_{i:j}$\\
Затримка & Низька & Вища, залежить від вікна\\
Обчислення & Рекурсивні матриці & Розріджена оптимізація\\
Корекція минулого & Зазвичай ні & Так\\
Типовий інструмент & KF/EKF/ESKF/UKF/PF & factor graph, bundle adjustment, fixed-lag smoother\\
\bottomrule
\end{longtable}

\section{Практичне правило}

Для вбудованого оцінювання стану в реальному часі першим перекладати треба фільтри: KF, EKF, ESKF, інновації, коваріації, шум. Для високоточних офлайн або віконних задач треба додати згладжування, факторні графи і розріджену лінійну алгебру.
